\documentclass[12pt]{article}
\usepackage[utf8]{inputenc}
\usepackage[T1]{fontenc}
\usepackage{hyperref}
\title{Western Frisian (Retired) Language Resources for K-12 Curriculum Development}
\author{Compiled by Claude Sonnet 5.5 (claude-sonnet-5-5) for Andr\'as Kornai}
\date{}
\begin{document}
\maketitle

\section{Introduction}
The code \texttt{fri} (Glottocode west2863) is not a language in its own right. Glottolog lists it as a retired ISO 639-3 entry (family field: ``Bookkeeping''), retired to align with ISO 639-2 in favour of \texttt{fry}, Western Frisian. Sources found in search: \url{https://glottolog.org/resource/languoid/id/west2863} and a 2008 IETF languages-list post noting that \texttt{fri} is ``retired'' in ISO 639-3: \url{https://www.alvestrand.no/pipermail/ietf-languages/2008-April/007752.html}.

\section{Primary Educational Resources}
None are attributable to \texttt{fri} as such. Educational materials for Western Frisian belong under the active code \texttt{fry} and were not searched here.

\section{Linguistic and Reference Materials}
Only the bookkeeping records cited above were found.

\section{Religious and Community Resources}
None found for this code.

\section{Digital and Audio Resources}
None found for this code.

\section{Curriculum Development Context}
A curriculum should be developed under \texttt{fry}. This dossier is deliberately short rather than padded.

\end{document}